\section{Lattice perturbation theory}
\label{sec:per}

In order to gain some insight into the non-perturbative results
presented in the next section, we have performed the corresponding
one-loop perturbative calculation both in the continuum and
on the lattice  in the Landau gauge. This calculation is an extension of
the one presented in
ref.\ \cite{newi}.

The standard calculation of the one-loop corrections
corresponds to a computation of  the renormalization constants with a fixed
ultra-violet
cut-off on the loop momenta, $\vert p_i \vert \le \pi/a$,  neglecting terms
of $O(a)$ in the final result. The renormalization constants obtained in
this limit depend on the quark mass $m_q$ and on the external momenta only
through the logarithmically divergent terms, which are related to the anomalous
dimension of the operator. Since, at one loop, these terms are universal, i.e.\
they are the same on the lattice and in any  continuum regularization,
the renormalization constants, necessary to relate lattice to continuum
renormalized operators, are independent of the quark mass and external momenta
(they depend however on the lattice spacing and on the continuum
renormalization scale $\mu$). We call this limit ``Standard Perturbation
Theory"
(STP).
The non-perturbative calculation however, is performed  on a lattice of finite
size, and at a fixed value of the lattice spacing and  quark
mass: thus we cannot make  terms of $O(a)$ arbitrarily small. In particular,
since we want to study the $\mu$-dependence of the renormalization constants,
it
is important to know at which values of $\mu$ (and $m_q$),
the perturbative results, obtained on
a finite lattice and with a fixed lattice spacing,
differ from those obtained in SPT. In  order to follow closely the
non-perturbative computation,   the perturbative  results have also been
obtained, on a lattice of size $16^3 \times 32$. In this way,
we keep
 all terms of $O(\alpha_s a)$ or $O(a^2)$, or higher, in the final
result (we call this procedure ``Discrete Perturbation Theory", DPT).   The
loop integrals were
computed  numerically, both in SPT and in DPT.

The non-perturbative results have been obtained by tracing suitable
amputated Green functions,  computed between off-shell quark states, in the
Landau gauge, cf.\ eq.(\ref{eq:proj}). In perturbation theory,  this
procedure differs from the standard one, where one identifies the
 correction from the term  proportional to the
tree-level matrix element of the operator under consideration \cite{buras,ciu}.
 The
difference is due to finite terms  which, in the standard procedure, are
considered as part of the matrix element of the operator, whereas  with the
projectors, are considered as part of the one-loop corrections. At
one-loop order,
these terms are the same in the continuum and on the lattice (in the limit $a
\to 0$) and cancel when one computes the difference between the continuum
and the lattice renormalization constants.
\begin{figure}[htbp]
\centering
\includegraphics[width=0.62\textwidth]{images/feynman2.pdf}
\caption{Vertex diagrams with the Clover action. The second and
third  diagrams
refer to the contribution from the rotated part of the operator.}
\label{fig:vertex}
\end{figure}

To be more specific, let us
consider the diagrams in fig.\ \ref{fig:vertex}, computed in
Na\"\i ve Dimensional Regularization (NDR), in the Feynman
gauge
\beq \Lambda_{\gamma_\mu}(\frac{p}{\mu})=
<p|\bar \psi \gamma_\mu \psi|p> =
 \Bigl[ 1+\frac{\alpha_s}{4 \pi} C_F \Bigl(\frac{p^2}{\mu^2}\Bigr)^{-\epsilon}
\Bigl( (\frac{1}{\hat \epsilon}+1) \gamma_\mu-2\frac{\pslash p_\mu}{p^2}
\Bigr) \Bigr]   ,\label{eq:gcont}\eeq
where $C_F=(N^2-1)/2N$, $1/\hat \epsilon=1/\epsilon
-\ln 4 \pi - \gamma_E$
and $p$ is the external momentum. Usually, one takes the
term proportional to $(1/\hat \epsilon +1)$, i.e.
 the coefficient of $\gamma_\mu$,
as the one loop contribution to the bare operator.
 With the projector (see eq.(\ref{eq:proj}))  the last term in
eq.\ (\ref{eq:gcont}) also gives a contribution, so that the factor now becomes
$(1/\hat \epsilon+1/2)$. In table \ref{tab:cont}, we give the
one-loop corrections obtained, with the projectors and in the
't Hooft and Veltman (HV) dimensional
renormalization scheme \cite{hv}, for the  two-quark operators considered in
this paper.
The NDR regularization differs from the HV one by the
definition of the anti-commuting properties of $\gamma_5$.
For details see for example \cite{ciu}.
 In NDR the corrections to the axial current and pseudoscalar density coincide
with those to the vector and scalar density respectively.
 By writing the gluon propagator as
$G(k)=\left(-\delta_{\mu\nu}k^2+(1-\lambda) k_\mu k_\nu \right)/k^4$,
 we denote as ``Landau" the contribution to the one-loop corrections
proportional to $(1-\lambda)$.
\begin{table}[htbp]
\begin{center}
\begin{tabular}{||c|c|c||}
%
\hline
%
\textrm{Operator}
&Feynman&Landau   \\ \hline \hline
%
%
$\gamma_\mu$ &$1/\hat \epsilon+ 1/2$ &$-1/\hat
\epsilon-1/2$  \\
\hline
$\gamma_\mu\gamma_5$&$1/\hat \epsilon+5/2$ &$-1/\hat
\epsilon-1/2$     \\ \hline
$\hat 1$  &$4/\hat \epsilon+6$ &$-1/\hat
\epsilon-2$ \\ \hline
%
$\gamma_5$  &$4/\hat \epsilon+10$ &$-1/\hat
\epsilon-2$ \\ \hline
%
$\Sigma$  &$-1/\hat \epsilon-1$ &$1/\hat
\epsilon+1$ \\  \hline
\end{tabular}
\caption[]{ \itshape The results of the calculation of the one-loop diagram of
fig.\ \ref{fig:vertex} for different two-quark operators, in the HV scheme, are
given. We also give the result of the calculation of the self-energy
diagram of fig.\ \ref{fig:self}.
 The mass of the quark has been taken to be
zero, the inverse quark propagator is written as $i \pslash \Sigma(p^2)$
and a factor $(\alpha_s/4 \pi) C_F \Bigl(p^2/\mu^2\Bigr)^{-\epsilon}$
 is omitted.}
\label{tab:cont}
\end{center}
\end{table}
\begin{figure}[htbp]
\centering
\includegraphics[width=0.62\textwidth]{images/feynman3.pdf}
\caption{Self energy diagram.}
\label{fig:self}
\end{figure}

With the $\overline{{\rm MS}}$ procedure, after the subtraction of the pole,
the
irreducible vertex of fig.\ \ref{fig:vertex} has the form\footnote{$
\Gamma_O(p/\mu)$ is obtained from $\Lambda_O(p/\mu)$ using a projector, as
described in eq.\ (\ref{eq:proj}).}
 \beq \Gamma_O(p/\mu)=
  1+\frac{\alpha_s}{4 \pi} C_F
\Bigl(-\gamma^{(O)} \ln (p^2/\mu^2)+C_O^{\overline{{\rm MS}}}
 \Bigr). \label{eq:gms}\eeq
{}From table \ref{tab:cont},  one  finds $C_{\gamma_\mu}^{\overline{{\rm
MS}}} =1/2$, $C_{\hat 1}^{\overline{{\rm MS}}}
 =6$, etc., in the Feynman gauge and
$C_{\gamma_\mu}^{\overline{{\rm MS}}} =-1/2$, $C_{\hat 1}^{\overline{{\rm MS}}}
 =-2$, etc., for the Landau (longitudinal) terms.

On the lattice, one gets a
similar expression \beq \Gamma_O(pa)=
  1+\frac{\alpha_s}{4 \pi} C_F
\Bigl(-\gamma^{(O)} \ln (p^2 a^2)+C_O^{{\rm Latt}}
 \Bigr). \label{eq:gl}\eeq
The values of $C_O^{{\rm Latt}}$, for the Feynman and Landau  corrections are
given in  table \ref{tab:latt}.
\begin{table}[htbp]
\begin{center}
\begin{tabular}{||c|c|c|c||}
%
\hline
%
\textrm{Operator}
&Feynman&Rotated Feynman&Landau   \\ \hline \hline
%
%
$\gamma_\mu$ &$6.62$ &$-3.52$ &$-4.29$ \\
\hline
$\gamma_\mu\gamma_5$&$5.09$ &
$-11.68$ &$-4.29$ \\ \hline
%
$\hat 1$ &$16.11$&$-12.00$&$-5.79$ \\ \hline
$\gamma_5$  &$19.18$ &$4.31$ &$-5.79$\\ \hline
%
$\Sigma$  &$8.21$ &$--$&$4.79$ \\  \hline
\end{tabular}
\caption[]{ \itshape The values of $C_O^{{\rm Latt}}$
 for different two-quark operators and for the self-energy
($C_\Sigma$), with the Clover
action. The mass of the quark has been taken to be
zero.}
\label{tab:latt}
\end{center}
\end{table}

In table \ref{tab:latt}, we have denoted by ``Rotated Feynman" the contribution
coming from the rotated part of the two-fermion operators, see eqs.
(\ref{eq:oi}) and (\ref{eq:rotaz}). The Landau contribution to the
rotated part of the operators vanishes. This is required by  the fact that the
renormalization constants which relate lattice to continuum operators are
gauge invariant.

Following refs.\ \cite{newi}--\cite{2ferm}, we write
\beq O(\mu)=\left(1+\frac{\as}{4 \pi} C_F \left(-\gamma^{(0)}\ln (\mu a)^2
+\Delta_O \right) \right) O(a) ,\eeq
where $\Delta_O$ can be computed from tables \ref{tab:cont} and \ref{tab:latt}
\beq \Delta_O=C_O^{\overline{{\rm MS}}}-C_O^{{\rm Latt}} +
 C_\Sigma^{\overline{{\rm MS}}}-C_\Sigma^{{\rm Latt}} .\eeq
In the difference between continuum and lattice one-loop corrections,
all  Landau pieces cancel, as expected. The results for
$\Delta_{\gamma_\mu},
\Delta_{\gamma_\mu \gamma_5},\cdots$ all agree with those of ref.
\cite{newi}.

On a lattice, for a given quark mass and
external momenta, there are terms of order $a$  which can become important
at large values of $m_q$ and $p$. Since with our method  we have to impose
the renormalization conditions at values of the momenta
which are large,
it is important to monitor, at least in perturbation theory, the values
of masses and momenta  at which $O(a)$ effects become important. The Green
function in DPT has a form similar to (\ref{eq:gl})
\beq \Gamma_O(pa)=
  1+\frac{\alpha_s}{4 \pi} C_F
\Bigl(-\gamma^{(O)} \ln (p^2 a^2)+C_O^{{\rm Latt}}(pa,m_q a)
 \Bigr), \label{eq:gld}\eeq
where, however, the ``constant" $C_O^{{\rm Latt}}$ now depends on the terms
of $O(a)$. Thus, in the calculation of the one-loop corrections in DPT,
 all the cut-off dependent terms have been included.
On a finite volume, the propagators of massless particles at zero momentum
must be regulated. The results of Feynman diagrams
depend on the regulator, however this dependence vanishes as the volume
is increased. In the results presented below, we have set the zero-momentum
terms in both the quark and gluon propagators to zero. We observe a small
difference in the results from SPT and DPT, at low values of $\mu^2$, which
is due to this choice. At the value of the quark mass at which we
have performed the simulation, mass corrections are negligible.

The coupling constant  appearing in eqs.\ (\ref{eq:gl}) and (\ref{eq:gld})
is the bare lattice coupling.
In ref.\ \cite{lepage}   it
was argued that it is possible to optimize the convergence of the perturbative
series by  reorganizing the  expansion in terms of a
``boosted"  coupling $\alpha_s^V$, defined in terms of a physical quantity.
In the comparison with the non-perturbative results below  we follow
the suggestion of ref.\ \cite{lepage}, and compute
the perturbative corrections using an effective coupling
\beq \alpha_s^V= \frac{1}{\langle \frac{1}{3}{\mathrm{Tr}} U_{plaq}\rangle}
 \alpha_s^{LATT} \simeq  1.68\,\, \alpha_s^{LATT} ,\eeq
where $\alpha_s^{LATT}=g_0^2/4 \pi$ and
$\langle \frac{1}{3}{\mathrm{Tr}} U_{plaq}\rangle$ is the expectation
value of the plaquette. We call this expansion ``Boosted Discrete
Perturbation Theory" (BDPT). Analogously we denote the  ``Boosted
Standard Perturbation Theory" by BSPT.

The results of the calculation of the
renormalization constants in BDPT will be compared to the non-perturbative
results in the next section.
